\subsubsection{品質属性}

設計品質に関わる属性には、モジュール性、保守容易性、移植性、テスト容易性、使いやすさ、正しさ、堅ろう性などがある。これらは、利害関係者の関心事の大きな部分を占める。

品質属性には、実行時に観察できるものと、実行時には直接見えにくいものがある。性能、セキュリティ、可用性、機能性、使いやすさは、実行時に観察しやすい。一方、変更容易性、移植性、再利用性、テスト容易性は、設計や保守の中で評価されることが多い。概念的一貫性、正しさ、完全性は、設計記述そのものを見て評価される。
